\section{Arithmetic closed points in every characteristic-zero principal open}

We isolate the arithmetic closed-fibre step used in spreading arguments.
The rings, principal opens, residue fields and quotient maps below are the
original ones. The finite-type Jacobson theorem and Zariski's lemma are the
standard results of \cite{Stacks}; the latter is used over $\mathbb Z$ itself.

\begin{lemma}[The integers are Jacobson]
\label{lem:arithmetic-integers-jacobson}
\lean{EnvelopeProperty.ArithmeticClosedPoints.int_isJacobsonRing}
\leanok
The ring $\mathbb Z$ is Jacobson.
\end{lemma}
\begin{proof}
Every nonzero prime ideal of the principal ideal domain $\mathbb Z$ is
maximal, and therefore equals its Jacobson radical. For the zero ideal,
let $a\ne0$ be an integer. Choose a prime $q>|a|$; the maximal ideal $(q)$
does not contain $a$, since $q\mid a$ would imply $q\le |a|$.
Thus the intersection of all maximal ideals contains no nonzero integer,
so it is zero. Consequently every prime ideal equals its Jacobson radical,
which is the Jacobson criterion.
\end{proof}

\begin{lemma}[The original closed residue fields are finite]
\label{lem:arithmetic-finite-closed-residue}
\lean{EnvelopeProperty.ArithmeticClosedPoints.quotient_prime_characteristic_and_finite}
\leanok
\uses{lem:arithmetic-integers-jacobson}
Let $R$ be a commutative algebra of finite type over $\mathbb Z$ and let
$\mathfrak m$ be any maximal ideal of $R$. There is a prime number $p$ such
that the original field $K=R/\mathfrak m$ has characteristic $p$, and $K$
is a finite set.
\end{lemma}
\begin{proof}
The original quotient is a finite-type $\mathbb Z$-algebra and is a field.
Zariski's lemma over a Jacobson ring makes $K$ a finite $\mathbb Z$-module;
in particular $\mathbb Z\to K$ is integral. If $K$ had characteristic zero,
that map would be injective. An injective integral ring map into a field
forces the source to be a field: for a nonzero source element, an integral
equation for its inverse in the field expresses that inverse as an element
of the source. This would make $\mathbb Z$ a field, a contradiction.
The nonzero characteristic of a field is prime; call it $p$. Its canonical
$\mathbb F_p$-algebra structure makes $K$ a finite-type $\mathbb F_p$-algebra:
the same finite family of integer-algebra generators also generates over
$\mathbb F_p$. Zariski's lemma over this field makes $K$ a finite
$\mathbb F_p$-module. A module generated by $d$ elements over the finite ring
$\mathbb F_p$ is the image of the finite set $\mathbb F_p^d$. Hence $K$ is
finite, as claimed.
\end{proof}

\begin{theorem}[Arbitrarily large characteristic closed points avoid the original element]
\label{thm:arithmetic-large-characteristic-closed-point}
\lean{EnvelopeProperty.ArithmeticClosedPoints.exists_large_characteristic_closed_point}
\leanok
\uses{lem:arithmetic-integers-jacobson,lem:arithmetic-finite-closed-residue}
Let $R$ be a characteristic-zero integral domain of finite type over
$\mathbb Z$. For every original nonzero $r\in R$ and every integer
$N\ge0$, there are an original maximal ideal $\mathfrak m$ of $R$ and a
prime $p>N$ such that $r\notin\mathfrak m$, the original quotient
$R/\mathfrak m$ has characteristic $p$, and $R/\mathfrak m$ is finite.
\end{theorem}
\begin{proof}
The integer $N!$ is nonzero in $R$ because $R$ has characteristic zero.
Thus $a=rN!$ is nonzero because $R$ is a domain. A finite-type algebra over
a Jacobson ring is Jacobson, so $R$ is Jacobson. Its zero ideal is radical,
and hence the intersection of all maximal ideals is zero. Some maximal
ideal $\mathfrak m$ therefore avoids $a$. Since $\mathfrak m$ is an ideal,
it avoids both $r$ and $N!$: membership of either factor would put their
product in $\mathfrak m$.

The preceding lemma supplies a prime characteristic $p$ and finiteness for
the original quotient $R/\mathfrak m$. If $p\le N$, then $p\mid N!$.
Since $p$ is zero in this quotient, the same is true of $N!$, contradicting
$N!\notin\mathfrak m$. Therefore $N<p$. All assertions concern this same
maximal ideal and its original quotient.
\end{proof}

\begin{corollary}[The arithmetic specialization is the original quotient map]
\label{cor:arithmetic-large-characteristic-quotient-map}
\lean{EnvelopeProperty.ArithmeticClosedPoints.exists_large_characteristic_quotient_map}
\leanok
\uses{thm:arithmetic-large-characteristic-closed-point}
Under the preceding hypotheses there are $\mathfrak m$ and $p>N$ such
that $K=R/\mathfrak m$ is finite of characteristic $p$, and the actual
quotient homomorphism $\varphi:R\to K$, $a\mapsto a\bmod\mathfrak m$,
is surjective and satisfies $\varphi(r)\ne0$.
\end{corollary}
\begin{proof}
Choose the maximal ideal and prime from the theorem. The original quotient
map is surjective by the definition of a quotient. Its kernel is exactly
$\mathfrak m$, so $r\notin\mathfrak m$ is precisely $\varphi(r)\ne0$.
\end{proof}

This statement supplies the arithmetic closed fibre in the principal open
$D(r)$. It asserts no descent of a family or line bundle, no lifting to
length-two Witt vectors, and no cohomological vanishing theorem.
